%% core/c2-branes.tex —— C2 轨道：A/B 模型与膜的类型、B-膜、A-膜、线算子、特征膜
%% 七段式：\oneline / definition / theorem+proof / \physside / example / \statusline / \dictionaryrow

\begin{centry}{A/B 模型与膜的类型}{A/B Models and Brane Types}{ab-model}
\oneline{同一个二维超对称 sigma 模型的两种拓扑扭化：A-模型只依赖复化 Kähler 类 $[B+i\omega]$，B-模型只依赖复结构 $\hat J$；靶空间超凯勒时，膜按三个复结构 $I,J,K$ 各标 A 型或 B 型，这一分类是整本词典的坐标系。}

\begin{definition}[A-模型与 B-模型]
设 $X$ 为实 $2n$ 维流形，带黎曼度量 $g$ 与闭二形式 $B\in\Omega^2(X)$（$B$ 场）。二维 sigma 模型的世界面为 $\Sigma$，玻色场为 $\Phi:\Sigma\to X$。拓扑扭化后未破缺的超对称方程要求一对复结构 $J_+,J_-$ 满足
\[
(1-iJ_+)\bar\partial\Phi=0,\qquad (1-iJ_-)\partial\Phi=0,
\]
即 $\Phi$ 关于 $J_+$ 全纯、关于 $J_-$ 反全纯。

若 $J_+=J_-=\hat J$，称为 \emph{B-模型}：$\Phi$ 关于 $\hat J$ 全纯，理论只依赖复结构 $\hat J$。若 $J_+=-J_-=\hat J$，称为 \emph{A-模型}：理论只依赖复化 Kähler 类 $[B+i\omega]$，其中 $\omega(\cdot,\cdot)=g(\hat J\,\cdot,\cdot)$。

若 $X$ 超凯勒，即带 $(g,I,J,K)$ 使 $\omega_I=gI$、$\omega_J=gJ$、$\omega_K=gK$ 均闭且 $K=IJ$，则二维扭化由参数 $t\in\mathbb{P}^1$ 标记：$t=\pm i$ 给出复结构 $\pm J$ 的 B-模型；实 $t$ 给出辛形式为 $\omega_I$ 与 $\omega_K$ 之实线性组合的 A-模型（$t=1$ 给 $K$，$t=0,\infty$ 给 $I$）；$t$ 与 $-t$ 给出等价的二维理论。
\end{definition}

\begin{definition}[膜的类型]
三类最基本的膜是：\emph{Lagrange 膜}，支撑 $L$ 满足 $\dim L=\frac12\dim X$ 且 $\omega|_L=0$；\emph{余各向同性膜}，支撑 $Y$ 满足 $TY^{\perp_\omega}\subseteq TY$（从而 $\dim Y\ge\frac12\dim X$）；\emph{全纯膜}，即 B-膜，支撑为复子流形并带全纯向量丛，更一般地是相干层（见 \ref{core:bbrane}）。

超凯勒靶上的膜是保持两个拓扑超荷的边界条件。这样的边界条件在复结构 $I,J,K$ 的每一个中要么是 B-型（保持该复结构的 B-模型超对称），要么是 A-型，于是带三格标签：$(B,B,B)$、$(B,A,A)$、$(A,B,A)$、$(A,A,B)$。典型地，$(B,B,B)$ 膜的支撑在三个复结构中全纯；$(B,A,A)$ 膜的支撑在 $I$ 中全纯、关于 $\omega_J$ 与 $\omega_K$ 为 Lagrange；$(A,B,A)$ 膜的支撑关于 $\omega_I,\omega_K$ 为 Lagrange、在 $J$ 中全纯。
\end{definition}

\begin{theorem}[{A-模型只看 $[B+i\omega]$}]
设 $X$ 紧 Kähler，$B,B'$ 为闭二形式，$\omega,\omega'$ 为 Kähler 形式，且 $[B'+i\omega']=[B+i\omega]\in H^2(X;\mathbb{C})$。则两个 A-模型的经典作用量相差一个超对称恰当项，从而给出同一个拓扑场论；而 B-模型只依赖复结构 $\hat J$，与 $\omega$ 无关。
\end{theorem}
\begin{proof}
A-模型的作用量可写成 $\{Q,V\}+\int_\Sigma\Phi^*(\omega-iB)$（Witten 1988 的构造），第二项只依赖上同调类 $[\omega-iB]$；而 $[B+i\omega]=i\,[\omega-iB]$，故同调假设给出相同的作用量。B-模型的作用量只含 $\hat J$ 与 $\bar\partial_{\hat J}$，$\omega$ 不出现；$B$ 的 $(1,1)$、$(2,0)$ 分量可经场重定义吸收，只有 $(0,2)$ 分量会扭化 B-膜范畴。
\end{proof}

\physside{把四维理论紧化在 $M=\Sigma\times C$ 上，低能有效理论就是靶空间为 $\mathcal M_H(G,C)$ 的二维拓扑 sigma 模型（KW06 §5 开头），而 $\mathcal M_H$ 超凯勒：复结构 $I$ 是「Higgs 丛 $(E,\varphi)$」的复结构，$J$ 是「平坦 $G_{\mathbb{C}}$-丛」的复结构，$K=IJ$ 是第三结构。数学侧「A-模型只依赖 $[B+i\omega]$」对应物理侧「A-模型看到的是辛形式 $\omega_K$ 与 $B$ 场的组合」；B-模型只依赖 $\hat J$ 对应「它只看到把 $\mathcal M_H$ 当成平坦丛模空间的那个复结构」。$S$-对偶在二维的化身是 $A(\mathcal M_H({}^L G),\omega_K)\simeq B(\mathcal M_H(G),J)$。}
\begin{example}[$G=SU(2)$、$g=2$ 的膜类型]
取 $C$ 为亏格 $2$ 曲线。$\dim_{\mathbb{C}}\mathrm{Bun}_{SU(2)}=3(g-1)=3$，故 $\dim_{\mathbb{C}}\mathcal M_H=\dim_{\mathbb{C}}T^*\mathrm{Bun}=6$，实维数 $12$。零膜支撑在一点，是 $(B,B,B)$ 型。Hitchin 纤维是复 $3$ 维环面（实维数 $6=\frac12\cdot12$），在 $I$ 中全纯、关于 $\omega_J,\omega_K$ 为 Lagrange，故带平坦 $U(1)$ 丛后是 $(B,A,A)$ 型。典范余各向同性膜支撑整个 $\mathcal M_H$，是 $(A,B,A)$ 型。
\end{example}

\statusline{定义与式 $(5.5)$--$(5.6)$、$[B+i\omega]$ 依赖：KW06 §5.1 p.65（式 $(5.33)$--$(5.34)$）。膜的四类型名单，以及不存在 $(B,B,A)$ 型与 $(A,A,A)$ 型膜：KW06 §5.5 pp.74--76 脚注（membrane type 总表 §5.6 p.77）。「B-膜范畴 $=D^b\mathrm{Coh}$」按物理推导处理，见 \ref{core:bbrane}。A-模型只依赖 $[B+i\omega]$ 是可逐项验证的数学事实（Witten 1988 的作用量构造）；超凯勒-复结构的对应见 Witten 2008 §2；膜类型的物理分类见 Frenkel 2009 §5.2。}

\dictionaryrow{B-膜与相干层：\ref{core:bbrane}；A-膜与 Lagrange 膜：\ref{core:abrane}；$\mathcal M_H$ 的超凯勒结构与三种复结构：\ref{core:hitchin}；扭化参数 $\Psi$ 如何取到 A/B 模型：\ref{core:gl-twist}、\ref{core:em-duality}。}
\end{centry}

\begin{centry}{B-膜＝相干层}{B-Branes as Coherent Sheaves}{bbrane}
\oneline{B-模型的开弦边界条件就是靶空间上的全纯对象：一条 B-膜由支撑子流形上的全纯向量丛（更一般地，相干层）给出，它们构成相干层的导出范畴 $D^b\mathrm{Coh}(X)$——词典里「膜 $=$ 层」的那一条边。}

\begin{definition}[B-膜]
设 $X$ 为复流形。一条 \emph{B-膜}是二元组 $(Y,F)$：闭复子流形 $Y\subseteq X$（可取 $Y=X$），全纯向量丛 $F\to Y$，带 $(0,1)$ 型联络 $\alpha$ 且 $\bar\partial_\alpha^2=0$（等价地，$F$ 带全纯结构）。允许把 $F$ 推广为支撑在 $Y$ 上的相干层；若 $B$ 场的 $(2,0)\oplus(0,2)$ 分量非零，则称相应的 $F$ 为扭化相干层。
\end{definition}

\begin{theorem}[{B-膜范畴（物理推导）}]
设 $X$ 紧致复流形，$B$ 场的 $(2,0)\oplus(0,2)$ 分量为零。则 B-模型保持超对称的开弦边界条件构成一个 $\mathbb{C}$-线性三角范畴，且与相干层的导出范畴等价：
\[
\mathrm{Brane}_B(X)\;\simeq\;D^b\mathrm{Coh}(X).
\]
对象级：$(Y,F)$ 对应支撑在 $Y$ 上的相干层 $\iota_*F$（$\iota:Y\hookrightarrow X$）。态射级：两条膜 $E,F$ 的开弦态空间是
\[
\operatorname{Hom}^{\bullet}(E,F)\cong\operatorname{Ext}^{\bullet}(E,F)
=\bigoplus_k H^k\bigl(X,\mathcal{E}xt^k(E,F)\bigr).
\]
\end{theorem}
\begin{proof}
边界上的费米子边界条件使 $\bar\partial_\alpha$ 成为世界面超荷的一部分；开弦 BRST 上同调于是按边界条件计算，恰是 $\bar\partial_\alpha$ 的上同调，即全纯截面，故边界条件由带全纯结构的丛/层分类。二维拓扑场论的开弦谱只有有限个非零次幂，第 $k$ 次正是 $H^k(X,\mathcal{H}om(E,F))$；开弦的粘接对应 Yoneda 复合，于是得到三角范畴结构。$B$ 的 $(0,2)$ 分量进入边界作用量的相位，等价于把 $\mathcal{O}_X$-模换成扭化 $\mathcal{O}_X$-模，故只改「扭化」而不改骨架。
\end{proof}

\physside{数学侧「$X$ 上的相干层 $E$」对应物理侧「B-模型在 $X$ 上的一条边界条件」；数学侧的 $\operatorname{Ext}^{\bullet}(E,F)$ 对应物理侧「两端贴 $E,F$ 的开弦的拓扑态空间」，而导出范畴中的平移与锥对应开弦态空间上的鬼数平移与开弦粘接——范畴的三角结构并不是额外数据，它就是开弦谱的物理结构。最重要的一条 B-膜是零膜 $\mathcal{O}_E$：支撑在一点 $x\in X$ 的天穹层；在 $\mathcal{M}_H({}^L G,C)$ 上，若 $x$ 对应平坦 ${}^L G_{\mathbb{C}}$-丛 $E$，则 $\mathcal{O}_E$ 是电特征膜（见 \ref{core:eigenbrane}）。$S$-对偶把零膜送到缠在 Hitchin 纤维上的 type $F$ 膜。B-膜与 A-膜的对偶两侧合起来就是镜像对称的三角。}
\begin{example}[$\mathbb{P}^1$ 与 $T^2$]
取 $X=\mathbb{P}^1$：$\operatorname{Ext}^{\bullet}(\mathcal{O},\mathcal{O})=H^{\bullet}(\mathbb{P}^1,\mathcal{O})=\mathbb{C}$（集中在 $0$ 次）；$\operatorname{Ext}^{\bullet}(\mathcal{O},\mathcal{O}(-1))=0$；$\operatorname{Ext}^{\bullet}(\mathcal{O},\mathcal{O}(-2))=H^{\bullet}(\mathbb{P}^1,\mathcal{O}(-2))=\mathbb{C}$（集中在 $1$ 次）。取 $X=T^2=\mathbb{C}/(\mathbb{Z}+i\mathbb{Z})$，零膜 $\mathcal{O}_x$：$\operatorname{Ext}^{\bullet}(\mathcal{O}_x,\mathcal{O}_x)=\wedge^{\bullet}T_x^*X$，次数 $0,1,2$ 的维数分别为 $1,2,1$，总维数 $4$。再取 $T^2$ 上次数 $0$ 的非平凡线丛 $M$：$h^0(M)=h^1(M)=0$，故 $\operatorname{Ext}^{\bullet}(\mathcal{O},M)=0$——不同的平坦线丛作为 B-膜互相正交。这些数直接由 Riemann--Roch 与 Koszul 复形算出。最后取 $X=T^*M$（$M$ 为紧流形）：$X$ 是复流形，其 B-膜就是相干层；而关于 Liouville 形式的 A-模型一侧，Lagrange 膜经 Nadler--Zaslow 对应到 $M$ 上可构造层的导出范畴——可看作 $M$ 的导出范畴在 Fourier--Mukai 变换下的像。这是「膜 $=$ 层」在 A-模型一侧的严格版本，细节待核。
\end{example}

\statusline{本条目按物理推导处理：「B-膜范畴 $=$ 相干层导出范畴」由 Kapustin 2006 论证（KW06 §5.3 p.70 采用；Frenkel 2009 §5.2 亦按此使用），但「B-模型的边界条件」本身没有数学定义，故不能算定理；对象级「B-膜 $=$ 相干层」的物理来源是 B-模型只依赖复结构（见 \ref{core:ab-model}）。可核对的数学事实是 $\operatorname{Ext}$ 群的计算（上例）。$T^*M$ 情形：$\mathcal{M}_H(G,C)$ 以 $T^*\mathcal{M}_{st}$ 为 Zariski 开集（Witten 2008 §3.1），而 $T^*M$ 的 Fukaya 范畴与 $M$ 上可构造层导出范畴的等价由 Nadler--Zaslow 2009 证明——这与本条目的「膜 $=$ 层」是同一类陈述，细节待核。}

\dictionaryrow{B-模型与膜类型：\ref{core:ab-model}；A-膜一侧：\ref{core:abrane}；零膜作为电特征膜：\ref{core:eigenbrane}；A-膜到 $D$-模的翻译器：\ref{core:dmod}、\ref{core:bcc}。}
\end{centry}

\begin{centry}{A-膜＝Lagrange 膜与余各向同性膜}{Lagrangian and Coisotropic Branes}{abrane}
\oneline{A-模型的开弦边界条件是 Lagrange 子流形上的平坦酉局部系统；$B$ 场开启时 Chan-Paton 曲率被要求与 $B$ 抵消。支撑维数高于中维时出现余各向同性膜，它是 A-膜到 $D$-模的翻译器的载体。}

\begin{definition}[A-膜]
设 $(X,\omega)$ 为实 $2n$ 维辛流形，$B$ 为闭二形式。一条 \emph{A-膜}是 $(L,\nabla)$：Lagrange 子流形 $L\subseteq X$（即 $\dim L=\frac12\dim X$ 且 $\omega|_L=0$），带 Hermite 线丛与酉联络 $\nabla$，其曲率 $F$ 满足
\[
F+B|_L=0 .
\]
$B=0$ 时这就是 $L$ 上的一个平坦酉局部系统，即 $\operatorname{Hom}(\pi_1(L),U(1))$ 中的元素。两条 Lagrange A-膜 $(L,\nabla)$ 与 $(L',\nabla')$ 之间的态空间是 Floer 上同调 $HF^{\bullet}((L,\nabla),(L',\nabla'))$；以 $(L,\nabla)$ 为对象、以 Floer 上同调为态射，就得到（带 $B$ 场扭化的）Fukaya 范畴。Chan-Paton 丛取直和得到多重膜，取秩 $N$ 的酉丛得到 $U(N)$ 情形。
\end{definition}

\begin{definition}[余各向同性膜]
子流形 $Y\subseteq X$ 称\emph{余各向同性}，若 $TY^{\perp_\omega}\subseteq TY$；等价地，$Y$ 局部是若干 Poisson 对易函数 $f_1,\dots,f_d$ 的公共零点集。于是 $\dim Y\ge\frac12\dim X$，等号时 $Y$ 是 Lagrange 的。当 $Y=X$、Chan-Paton 为线丛且曲率为 $F$ 时，A-膜条件为
\[
(\omega^{-1}F)^2=-1,
\]
此时 $\omega-iF$ 是 $X$ 上的全纯辛形式，故要求 $\omega$ 能写成某个全纯辛形式的实部。当 $Y\neq X$ 时，$B$ 场以 $F+B|_Y$ 进入同一条件：余各向同性 A-膜的存在性受 $[B|_Y]$ 的整性障碍控制，且 $B$ 的 $(2,0)\oplus(0,2)$ 分量会改变膜的存在性与模空间。
\end{definition}

\begin{theorem}[Lagrange A-膜的存在与分类]
设 $L\subseteq X$ 为紧 Lagrange 子流形，$B$ 为闭二形式。则：
\begin{enumerate}[label=(\roman*)]
\item $L$ 上存在秩 $1$ 的 A-膜，当且仅当 $[B|_L]\in2\pi H^2(L;\mathbb{Z})$；
\item 存在时，$L$ 上 A-膜的同构类构成 $\operatorname{Hom}(\pi_1(L),U(1))$ 的一个 torsor。
\end{enumerate}
\end{theorem}
\begin{proof}
A-膜条件要求 $F=-B|_L$，而 $F$ 是某个 $U(1)$ 联络的曲率，故 $[F]/2\pi\in H^2(L;\mathbb{Z})$ 是必要条件。反之，若 $[B|_L]$ 的周期是 $2\pi$ 的整数倍，则存在曲率为 $-B|_L$ 的 $U(1)$ 联络。两个这样的联络之差是闭 $1$-形式；模去规范等价（恰当形式）即模去 $H^1(L;\mathbb{Z})$，于是同构类构成 $H^1(L;\mathbb{R})/H^1(L;\mathbb{Z})\cong\operatorname{Hom}(\pi_1(L),U(1))$ 的 torsor。
\end{proof}

\physside{在二维 sigma 模型中，A-膜就是保持 A-扭曲超对称（$J_+=-J_-$）的边界条件；$B$ 场进入边界作用量的方式正是把 Chan-Paton 曲率平移 $F\mapsto F+B|_L$，所以 $F+B|_L=0$ 就是 A-膜条件（KW06 §5.2）。取靶为 $\mathcal M_H(G,C)$ 的 A-模型（辛形式 $\omega_K$）：Hitchin 纤维是复子流形、关于 $\omega_J,\omega_K$ 为 Lagrange，带秩 $1$ 平坦酉 Chan-Paton 丛后就是 $(B,A,A)$ 型膜（type $F$ 膜），即磁特征膜。典范余各向同性膜 $B_{cc}$ 取 $F=(\operatorname{Im}\tau)\omega_J$，是 $(A,B,A)$ 型，它把 A-膜翻译成 $D$-模。}
\begin{example}[$T^2$ 与 Hitchin 纤维]
取 $X=T^2$ 带平面积形式：一维 Lagrange 子流形是斜率 $(p,q)$（$\gcd(p,q)=1$）的闭测地线，其长度平方为 $p^2+q^2$（取单位面积、$\tau=i$）；曲率为 $0$ 的平凡线丛给出其上的 A-膜，模空间为 $\operatorname{Hom}(\mathbb{Z},U(1))\cong U(1)$。因 $\dim L=1$，$B|_L=0$ 自动成立，任何 $B$ 都不改变存在性。取 $X=\mathbb{P}^1$ 带 Fubini--Study 形式：$X$ 是复曲线，任何嵌入圆都是 Lagrange 的（$\omega$ 对一维子空间自动为零），如赤道 $S^1$；点是 B-膜而非 A-膜（$\dim 0\neq\frac12\dim 2$），这正显示 A、B 两类膜在维数上的差别。取 $G=SU(2)$、$g=2$：Hitchin 纤维是实 $6$ 维环面 $T^6$，恰等于 $\frac12\dim_{\mathbb{R}}\mathcal M_H$，带平坦 $U(1)$ 丛后是秩 $1$ 的 $(B,A,A)$ 型 A-膜；这类膜的模空间是参平坦 $U(1)$ 丛的复环面，其复维数等于纤维的复维数 $3$。
\end{example}

\statusline{定义 $F+B|_L=0$ 与 A-膜：KW06 §5.2 p.69；余各向同性条件与 $(\omega^{-1}F)^2=-1$：KW06 §11.1 p.185（式 $(11.1)$--$(11.2)$）。定理 (i)(ii) 是本词典给出的数学陈述（联络论与 de Rham 上同调的初等推论，证明见上），KW06 未以定理形式陈述。物理推导：余各向同性 A-膜的存在性由 Kapustin--Orlov 2003 论证；c.c. 膜 $B_{cc}$ 是 $(A,B,A)$ 型见 KW06 §11.1 pp.185--187。「A-膜范畴 $=$ Fukaya 范畴」在数学上有严格定义（Fukaya 等），但它与物理 A-模型边界条件的等同是猜想（Frenkel 2009 §5.2--5.3）。}

\dictionaryrow{A-模型与 B-模型：\ref{core:ab-model}；B-膜一侧：\ref{core:bbrane}；Hitchin 纤维作为复合 Lagrange 膜：\ref{core:hitchin-fibration}；典范余各向同性膜作为翻译器：\ref{core:bcc}。}
\end{centry}

\begin{centry}{线算子：Wilson 与 't Hooft}{Line Operators: Wilson and 't Hooft}{line-operators}
\oneline{沿一条曲线的非局部算子：Wilson 线取联络沿圈的平行移动之迹，'t Hooft 线靠指定沿圈的奇异性（磁通量）定义；在拓扑扭化理论中两者都是保超对称的算子，并在 $S$-对偶下互换。}

\begin{definition}[Wilson 线与 't Hooft 线]
设 $G$ 紧连通李群，$M$ 黎曼流形，$A$ 为 $G$-联络。对嵌入的定向曲线 $C\subseteq M$ 与 $G$ 的有限维酉表示 $R$，定义 \emph{Wilson 线}
\[
W_R(C)=\operatorname{Tr}_R\operatorname{Hol}_C(A)=\operatorname{Tr}_R P\exp\Bigl(\oint_C A\Bigr);
\]
若 $C$ 有端点 $p,q$，则不取迹而取平行移动矩阵 $\operatorname{Hol}_C(A)\in\operatorname{Hom}(P_p,P_q)$。

\emph{'t Hooft 线} $T(\lambda,C)$ 不要求联络光滑，而是指定沿 $C$ 的奇异性：取同态 $\rho:U(1)\to G$ 的共轭类，要求 $C$ 的管状邻域 $C\times D^3$ 内的场在横截切片 $D^3$（边界为 $S^2$）中具有以 $\rho$ 为单极子嵌入的奇异性。$G=U(1)$ 时即 Dirac 单极子型：$F=(i/2)\star d(1/|x|)$，$\int_{S^2}iF/2\pi=1$；一般情形写 $A=\rho(A_0)+\hat A$，要求 $\hat A$ 在 $C$ 上光滑且与 $\rho(A_0)$ 对易。
\end{definition}

\begin{proposition}[GNO 分类]
设 $G$ 紧半单，$C\subseteq M$ 为嵌入曲线且其法丛定向。则 't Hooft 线的（拓扑）荷由 ${}^L G$ 的不可约表示 $\hat R$ 分类；等价地，由同态 $\rho:U(1)\to G$ 的共轭类（即 $G$ 的支配余权，等价于 ${}^L G$ 的支配权）分类。
\end{proposition}
\begin{proof}
$C$ 的管状邻域同胚于 $C\times D^3$，沿 $C$ 的奇异性由横截切片 $S^2=\partial D^3$ 上 $G$-丛的拓扑类与局部模型共同决定。$S^2$ 上的 $G$-丛由 $\pi_1(G)$ 分类（Grothendieck），而 $\pi_1(G)\cong Z({}^L G)$，恰是 ${}^L G$ 的表示在中心上的特征标。局部模型由 $\rho$ 的共轭类给定，其最高权给出 ${}^L G$ 的支配权，两者合起来正是 ${}^L G$ 的一个不可约表示。反之，给定 ${}^L G$ 的不可约表示可构造相应奇异性。
\end{proof}

\physside{在四维 $\mathcal{N}=4$ 理论的拓扑扭化里，Wilson 线是有序算子、't Hooft 线是无序算子。$t=v/u=\pm i$ 时 Wilson 线用复化联络 $C=A+i\varphi$：$W(R,C)=\operatorname{Tr}_R P\exp\bigl(-\oint_C(A+i\varphi)\bigr)$（书 (6.4)）；$t\neq\pm i$ 时 holonomy 不再是拓扑不变量。$S$-对偶交换 Wilson 与 't Hooft，同时交换 $G$ 与 ${}^L G$。线算子的 OPE 给出算子代数：$W_R(C)W_{\tilde R}(C')=\sum_\alpha n_\alpha W_{R_\alpha}(C')$（书 (6.27)），系数 $n_\alpha$ 由表示分解 $R\otimes\tilde R=\bigoplus_\alpha n_\alpha R_\alpha$ 给出，两点函数给出 Zamolodchikov 型度量。更一般地，可同时带电、磁荷的混合算子 $WT_{n,m}$ 只在 $\Psi=-n/m$ 为有理数时存在：磁荷为 $m$ 的 't Hooft 奇异性背景下，规范不变性要求 Wilson 荷满足 $n=-m\Psi$ 为整数（书 (6.14)--(6.17)）；在拓扑扭化理论中全部线算子都由 $\Psi=\infty$ 的 Wilson 线经 $S$-对偶生成，故它们的荷成比例。把静态线算子沿紧化的时间圈取迹，得到相应表示中态空间的维数，这是线算子与边界态空间 $H_{B_1,B_2}$ 的联系。}
\begin{example}[$G=U(1)$ 与 $G=SU(2)$]
$G=U(1)$：$W_n(C)=\exp(in\oint_C A)$，$n\in\mathbb{Z}$；'t Hooft 线 $T_m$ 由 $m\in\mathbb{Z}$ 单位磁通量刻画；电荷对 $(n,m)\in\mathbb{Z}^2$，$S$-对偶作用为 $(n,m)\mapsto(-m,n)$。$G=SU(2)$：Wilson 线由自旋 $j\in\frac12\mathbb{Z}_{\ge0}$ 标记（即 $SU(2)$ 的不可约表示）；'t Hooft 线由 ${}^L G=SO(3)$ 的不可约表示标记，即整数自旋 $k\in\mathbb{Z}_{\ge0}$，维数 $1,3,5,\dots$——半整数自旋的 Wilson 线没有对偶的 't Hooft 线。对 $G=U(N)$ 有 ${}^L G=U(N)$（自对偶），故 't Hooft 线的荷由 $U(N)$ 的不可约表示分类，与 Wilson 线同型；这是书 §6.2 p.84 讨论的 GNO 对偶的自对偶情形。
\end{example}

\statusline{Wilson 线定义：KW06 §6.1 p.79（式 $(6.1)$）；'t Hooft 线与 GNO 分类：KW06 §6.2 pp.81--84（p.84 属 §6.2，§7 起于 p.103）。't Hooft 线的分类（Goddard--Nuyts--Olive 1977）是半经典物理推导，Kapustin 2006 重释为算子对偶；$S^2$ 上 $G$-丛由 $\pi_1(G)$ 分类是 Grothendieck 的定理。OPE 系数 $(6.27)$ 在超对称下可由表示论经典计算；'t Hooft 侧 OPE $(6.29)$ 只是预期（书中说物理文献未直接算过）。Wilson 线的拓扑不变性只在 $t=\pm i$（$\Psi=\infty$）成立。Wilson 线算子在几何 Langlands 中的范畴化（Wilson 函子 $W_V$）见 Frenkel 2009 §3。}

\dictionaryrow{$S$-对偶交换 Wilson 与 't Hooft：\ref{core:em-duality}；线算子作用的边界条件即特征膜：\ref{core:eigenbrane}；'t Hooft 算子等于 Hecke 算子：\ref{core:hecke}、\ref{core:monopole-hecke}；拓扑扭化的参数 $\Psi$：\ref{core:gl-twist}。}
\end{centry}

\begin{centry}{特征膜}{Eigenbranes}{eigenbrane}
\oneline{特征膜是「线算子作用在边界条件上只贡献一个向量空间因子」的膜：$XB\simeq B\otimes V$。电特征膜是 Wilson 线的共同特征膜，磁特征膜是 't Hooft 线的共同特征膜，其 B-模型侧就是 Beilinson--Drinfeld 的 Hecke 特征层。}

\begin{definition}[特征膜]
设 $\mathcal{T}$ 为二维拓扑场论，$B$ 为其一个边界条件（膜）；设 $X$ 为一条线算子，支撑在平行于边界的直线 $L$ 上。把 $L$ 推向边界，$X$ 把 $B$ 变成新膜 $XB$。若存在有限维向量空间 $V$ 与同构
\[
XB\simeq B\otimes V,
\]
则称 $B$ 是 $X$ 的\emph{特征膜}，$V$ 为本征值。对一族交换的线算子 $\{X_i\}$，共同特征膜要求对每个 $i$ 有 $X_iB\simeq B\otimes V_i$。若线算子 $X_p$ 对 $C$ 上每点 $p$ 都有定义，则 $V_p$ 随 $p$ 构成 $C$ 上的平坦向量丛（可带非平凡单值），这给出「特征膜 $=$ 对象在 $C$ 上的平坦丛」的直观。特征膜在数学侧的化身是 Beilinson--Drinfeld 的 Hecke 特征层：Hecke 算子 $H_R$ 的特征层 $F$ 满足 $H_R(F)\cong V_R\otimes F$，正是 B-膜意义上的特征对象（见 \ref{core:hecke}）。

\emph{电特征膜}：$\Psi=\infty$（纯 Wilson 扭化）处 Wilson 线的共同特征膜；\emph{磁特征膜}：$\Psi=0$（纯 't Hooft 扭化）处 't Hooft 线的共同特征膜。$S$-对偶交换两者。在 §5.5 的四类型名单里，磁特征膜（type $F$ 膜）是 $(B,A,A)$ 型；电特征膜则可由零膜这类 $(B,B,B)$ 型膜的扭化直和实现（\S8.2）。电、磁特征膜的本征值空间与离散电、磁通量 $e_0,m_0$ 的平移相容（书 (8.8)--(8.9)），故 §7 的离散通量分类决定了特征膜的扭化类。
\end{definition}

\begin{theorem}[{零膜是电特征膜（物理推导）}]
设 $G$ 紧单连通，$C$ 紧黎曼面，$\Psi=\infty$。设 $x\in\mathcal M_H({}^L G,C)$ 为光滑点（对应既约不可约的平坦 ${}^L G_{\mathbb{C}}$-丛），$B_x$ 为支撑在 $x$ 的零膜。则对每个 $p\in C$ 与 ${}^L G$ 的表示 $R$，
\[
X_p(R)\,B_x\simeq B_x\otimes E_p(R)|_x,
\]
其中 $E(R)\to\mathcal M_H({}^L G)\times C$ 是泛 Higgs 丛的第 $R$ 部分，$E_p(R)=E(R)|_{\mathcal M_H\times\{p\}}$。故 $B_x$ 是电特征膜，本征值 $V_p=E_p(R)|_x$ 构成 $C$ 上的平坦丛。
\end{theorem}
\begin{proof}
Wilson 线 $X_p(R)$ 作用于膜的方式是把 Chan-Paton 丛 $U$ 换成 $U\otimes E_p(R)$（书 (8.7)）。零膜的 Chan-Paton 层是点 $x$ 处的天穹层 $\delta_x$，而 $\delta_x\otimes E_p(R)\simeq\delta_x\otimes E_p(R)|_x$（书 (8.10)）——张量积只在 $x$ 处的纤维非零，故作用只贡献一个向量空间因子。$E(R)$ 关于 $\mathcal M_H$ 的每个复结构全纯，故该操作保持 $(B,B,B)$ 型超对称。$p$ 变动时 $E_p(R)|_x$ 是 $C$ 上平坦丛的纤维。
\end{proof}

\physside{几何 Langlands 的两侧在此对译：数学侧「平坦 ${}^L G$-丛 $E$ 对应 Hecke 特征 $D$-模 $F_E$」对应物理侧「电特征膜 $B_{x(\vartheta)}$ 对应磁特征膜」。$S$-对偶把 Wilson 与 't Hooft 互换，于是把电特征膜送到磁特征膜：支撑在 Hitchin 纤维 $F$ 上、带秩 $1$ 平坦酉 Chan-Paton 的 type $F$ 膜。磁特征膜在 $\mathcal M_H(G,C)$ 上，再经典范余各向同性膜 $B_{cc}$ 的翻译器变成 ${\rm Bun}_G$ 上的 $D$-模，其 Hecke 特征值由 $E$ 的表示丛给出。用一句话说：特征膜就是「Wilson 线在该膜上的作用等于标量」的膜。}
\begin{example}[$G=U(1)$]
取 $G=U(1)$、$C$ 为椭圆曲线，则 ${}^L G=U(1)$，Wilson 线 $W_n$ 与 't Hooft 线 $T_m$ 的特征值分别为 $n,m\in\mathbb{Z}$。电特征膜是支撑在光滑点的零膜（本征值取该点处平坦 $U(1)$-丛的纤维）；磁特征膜是支撑在 Hitchin 纤维 $\check T^2$ 上带平坦线丛的膜。$S$-对偶把两种膜互换，并把纤维上的点送到对偶纤维上的对偶点，这与 $\mathcal M_H(U(1),C)$ 的 Hitchin 纤维互为对偶环面相容。
\end{example}

\statusline{定义 (6.22)、(6.24)、(6.25)：KW06 §6.4 p.95（物理定义，书中明言对应 Beilinson--Drinfeld 的 Hecke 特征层）。电特征膜与磁特征膜：KW06 §6.4 p.97。定理（零膜是电特征膜）是 KW06 §8.2 的物理推导（式 (8.7)、(8.10)），不是数学定理；「磁特征膜 $=$ type $F$ 膜」由 $S$-对偶论证（KW06 §9 p.125），其与单极子、Hecke 修正空间的联系是 KW06 §10 的主题（'t Hooft 算子即 Hecke 算子）。离散电、磁通量的分类见 KW06 §7 pp.103--104。「电特征膜对应平坦 ${}^L G$-丛」的数学对应就是几何 Langlands：$E$ 不可约时 Hecke 特征 $D$-模的存在性已有定理，但完整范畴等价仍是猜想（Frenkel 2009 §2--§3）。}

\dictionaryrow{线算子与特征值的来源：\ref{core:line-operators}；'t Hooft 算子等于 Hecke 算子：\ref{core:hecke}；几何 Langlands 的物理化身：\ref{core:gl}；$D$-模与平坦联络：\ref{core:dmod}。}
\end{centry}
